% --- IMPORTS ---
\documentclass[leqno]{article}
\usepackage{graphicx}
\usepackage{dsfont}
\usepackage{mathtools}
\usepackage{amssymb}
\usepackage{amsmath}
\usepackage{amsthm}
\usepackage{float}
\usepackage{ragged2e}
\usepackage{tensor}
\usepackage{fancyhdr}
\usepackage{empheq}
\usepackage[most]{tcolorbox}
\usepackage{enumitem}
\usepackage{dirtytalk}
\usepackage{marginnote}
\usepackage{microtype}
\usepackage[
  a4paper,
  left=0.8in,
  right=1in,
  top=1in,
  bottom=0.5in,
  headheight=14pt,
  headsep=0.4in
]{geometry}
\setlength{\parindent}{0cm}
% --- TikZ Packages ---
\usepackage{pgfplots}
\pgfplotsset{compat=1.18}
\usepgfplotslibrary{fillbetween, polar}
\usetikzlibrary{calc, positioning}
\usetikzlibrary{angles, quotes}
\usetikzlibrary{arrows.meta}
\usetikzlibrary{decorations.pathreplacing, decorations.markings}
\usetikzlibrary{patterns}
\usetikzlibrary{intersections}
% --- URL Packages ---
\usepackage[colorlinks=true,linkcolor=red,citecolor=magenta,urlcolor=black]{hyperref}%
\urlstyle{same}
% --- END OF IMPORTS ---

% --- CUSTOM COMMANDS ---
\RenewDocumentCommand{\marks}{o m}{%
  \IfValueTF{#1}%
    {\marginnote{\raggedright\textbf{[#2]}}[#1]}%
    {\marginnote{\raggedright\textbf{[#2]}}}%
}
\newcommand{\vspacerule}[2]{%
  \par\vspace*{#1}%
  \noindent\hrulefill\par
  \vspace*{#2}%
}
\definecolor{scarlet}{RGB}{205, 40, 75}
\newcommand{\questionitem}{%
  \item\phantomsection
  \addcontentsline{toc}{section}{Question \number\value{enumi}}%
}
\renewcommand{\contentsname}{Questions}
% --- END OF CUSTOM COMMANDS ---

% --- HEADER CONFIG ---
\pagestyle{fancy}
\fancyhead{}
\fancyfoot{}
\renewcommand{\headrulewidth}{0.9pt}
\lhead{\textit{Solutions} - Linear Transformations}
\chead{}
\rhead{\href{https://danieljsmith.org}{danieljsmith.org}}
% --- END OF HEADER CONFIG ---

\begin{document}
\vspace*{20pt}
\tableofcontents
\newpage
\begin{enumerate}[
  label=\textbf{\arabic*.},
  leftmargin=*,
  topsep=0pt,
  partopsep=0pt,
  parsep=0pt
]

% ---- Question 1 ----
\questionitem

\begin{enumerate}[label=\textbf{(\roman*)}]
  \item In each of the following cases, find a $2\times 2$ matrix that represents
  \begin{enumerate}[label=(\alph*)]
    \item a reflection in the line $y=x$ \marks{1}
    \item a rotation of $60^\circ$ clockwise about $(0,0)$ \marks{1}
    \item a rotation of $60^\circ$ clockwise about $(0,0)$ followed by a reflection in the line $y=x$ \marks{2} \\
  \end{enumerate}
  \item The triangle $T$ has vertices at the points $(1,2)$, $(5,5)$ and $(k,-1)$, where $k>-3$ is a constant. \\

  Triangle $T$ is transformed onto the triangle $T'$ by the matrix
  \[
  \begin{pmatrix}
  2 & 5 \\
  1 & -1
  \end{pmatrix}
  \]
  Given that the area of triangle $T'$ is $84$ square units, find the value of $k$. \marks{6} \\
\end{enumerate}

\vspace*{30pt}

\begin{tcolorbox}[colback=white, colframe=scarlet, title={\textbf{Solution}}, breakable, grow to left by=6mm, grow to right by=10mm]
\begin{enumerate}[label=\textbf{(\roman*)}]
  \item For a $2\times 2$ transformation matrix, the first column is the image of $\begin{pmatrix}1\\0\end{pmatrix}$ and the second column is the image of $\begin{pmatrix}0\\1\end{pmatrix}$.
  
  \begin{enumerate}[label=\textbf{(\alph*)}]
    \item A reflection in the line $y=x$ swaps the coordinates, so
    \[
    \begin{pmatrix}1\\0\end{pmatrix}\mapsto \begin{pmatrix}0\\1\end{pmatrix},
    \qquad
    \begin{pmatrix}0\\1\end{pmatrix}\mapsto \begin{pmatrix}1\\0\end{pmatrix}
    \]
    Hence the matrix is
    \[
    \begin{pmatrix}
    0 & 1\\
    1 & 0
    \end{pmatrix}
    \]

    \item A clockwise rotation through angle $\theta$ about the origin has matrix
    \[
    \begin{pmatrix}
    \cos\theta & \sin\theta\\
    -\sin\theta & \cos\theta
    \end{pmatrix}
    \]
    With $\theta=60^\circ$,
    \[
    \cos 60^\circ=\frac12,
    \qquad
    \sin 60^\circ=\frac{\sqrt3}{2}
    \]
    So the matrix is
    \[
    \begin{pmatrix}
    \frac12 & \frac{\sqrt3}{2}\\
    -\frac{\sqrt3}{2} & \frac12
    \end{pmatrix}
    \]

    \item Let
    \[
    R=
    \begin{pmatrix}
    \frac12 & \frac{\sqrt3}{2}\\
    -\frac{\sqrt3}{2} & \frac12
    \end{pmatrix},
    \qquad
    F=
    \begin{pmatrix}
    0 & 1\\
    1 & 0
    \end{pmatrix}
    \]
    Since the rotation is done first and the reflection second, the required matrix is $FR$.
    \begin{align*}
    FR&=
    \begin{pmatrix}
    0 & 1\\
    1 & 0
    \end{pmatrix}
    \begin{pmatrix}
    \frac12 & \frac{\sqrt3}{2}\\
    -\frac{\sqrt3}{2} & \frac12
    \end{pmatrix}\\
    &=
    \begin{pmatrix}
    0\cdot \frac12+1\cdot\left(-\frac{\sqrt3}{2}\right) &
    0\cdot \frac{\sqrt3}{2}+1\cdot \frac12\\[4pt]
    1\cdot \frac12+0\cdot\left(-\frac{\sqrt3}{2}\right) &
    1\cdot \frac{\sqrt3}{2}+0\cdot \frac12
    \end{pmatrix}\\
    &=
    \begin{pmatrix}
    -\frac{\sqrt3}{2} & \frac12\\
    \frac12 & \frac{\sqrt3}{2}
    \end{pmatrix}
    \end{align*}
    So the matrix is
    \[
    \begin{pmatrix}
    -\frac{\sqrt3}{2} & \frac12\\
    \frac12 & \frac{\sqrt3}{2}
    \end{pmatrix}
    \]
  \end{enumerate}

  \item Let
  \[
  A=
  \begin{pmatrix}
  2 & 5\\
  1 & -1
  \end{pmatrix}
  \]
  The determinant of $A$ is
  \[
  \det(A)=2(-1)-5(1)=-7
  \]
  A transformation by a matrix scales area by the factor $|\det(A)|$, so the area scale factor is
  \[
  |-7|=7
  \]
  Hence the area of triangle $T$ is
  \[
  \frac{84}{7}=12
  \]

  Now find the area of $T$ in terms of $k$.

  Using the vertex $(1,2)$, the side vectors are
  \[
  \begin{pmatrix}5-1\\5-2\end{pmatrix}
  =
  \begin{pmatrix}4\\3\end{pmatrix},
  \qquad
  \begin{pmatrix}k-1\\-1-2\end{pmatrix}
  =
  \begin{pmatrix}k-1\\-3\end{pmatrix}
  \]
  So
  \begin{align*}
  \text{Area of }T
  &=\frac12\left|\det
  \begin{pmatrix}
  4 & k-1\\
  3 & -3
  \end{pmatrix}
  \right|\\
  &=\frac12|4(-3)-3(k-1)|\\
  &=\frac12|-12-3k+3|\\
  &=\frac12|-3k-9|\\
  &=\frac32|k+3|
  \end{align*}

  Since $k>-3$, we have $k+3>0$, so
  \[
  |k+3|=k+3
  \]
  Therefore
  \[
  \frac32(k+3)=12
  \]
  \[
  k+3=8
  \]
  \[
  k=5
  \]

  Hence the value of $k$ is $5$.
\end{enumerate}
\end{tcolorbox}


\newpage

% ---- Question 2 ----
\questionitem

\[
\mathbf{P}=\begin{pmatrix} p+1 & 2 \\ 2 & p+1 \end{pmatrix}
\] \\
where $p$ is a positive constant. \\

The matrix $\mathbf{P}$ represents a linear transformation $U$. \\
The parallelogram $S$ has vertices at the points with coordinates $(1,1)$, $(4,2)$, $(6,5)$ and $(3,4)$. \\
The area of the image of $S$ under the linear transformation $U$ is $35$. \\

\begin{enumerate}[label=\textbf{(\alph*)}]
  \item Determine the value of $p$. \marks{4} \\
\end{enumerate}

The transformation $V$ consists of a reflection in the line $y=x$ followed by a stretch of scale factor $2$ parallel to the $x$-axis with the $y$-axis invariant. The transformation $V$ is represented by the matrix $\mathbf{Q}$. \\

\begin{enumerate}[label=\textbf{(\alph*)}, resume]
  \item Write down the matrix $\mathbf{Q}$. \marks{2} \\
\end{enumerate}

Given that $U$ followed by $V$ is the transformation $W$, which is represented by the matrix $\mathbf{R}$ \\

\begin{enumerate}[label=\textbf{(\alph*)}, resume]
  \item find the matrix $\mathbf{R}$. \marks{2} \\
\end{enumerate}

\vspace*{30pt}

\begin{tcolorbox}[colback=white, colframe=scarlet, title={\textbf{Solution}}, breakable, grow to left by=6mm, grow to right by=10mm]
\begin{enumerate}[label=\textbf{(\alph*)}]
  \item Take two adjacent sides of the parallelogram from the vertex $(1,1)$:
\[
\begin{pmatrix}4\\2\end{pmatrix}-\begin{pmatrix}1\\1\end{pmatrix}
=
\begin{pmatrix}3\\1\end{pmatrix},
\qquad
\begin{pmatrix}3\\4\end{pmatrix}-\begin{pmatrix}1\\1\end{pmatrix}
=
\begin{pmatrix}2\\3\end{pmatrix}
\]

So the area of $S$ is
\[
\left|\det\begin{pmatrix}3&2\\1&3\end{pmatrix}\right|
=|3\cdot 3-2\cdot 1|
=|9-2|
=7
\]

A linear transformation represented by matrix $\mathbf P$ multiplies area by the factor $|\det \mathbf P|$.

Since the image of $S$ has area $35$,
\[
35=|\det \mathbf P|\times 7
\]
so
\[
|\det \mathbf P|=5
\]

Now
\[
\mathbf P=\begin{pmatrix}p+1&2\\2&p+1\end{pmatrix}
\]
so
\begin{align*}
\det \mathbf P
&=(p+1)(p+1)-2\cdot 2\\
&=(p+1)^2-4\\
&=p^2+2p+1-4\\
&=p^2+2p-3
\end{align*}

Hence
\[
|p^2+2p-3|=5
\]

So either
\[
p^2+2p-3=5
\]
or
\[
p^2+2p-3=-5
\]

From the first:
\begin{align*}
p^2+2p-8&=0\\
(p+4)(p-2)&=0
\end{align*}
so
\[
p=-4 \quad \text{or} \quad p=2
\]

From the second:
\[
p^2+2p+2=0
\]
which has no real solutions.

Since $p$ is positive,
\[
p=2
\]

  \item Reflection in the line $y=x$ has matrix
\[
\begin{pmatrix}0&1\\1&0\end{pmatrix}
\]

A stretch of scale factor $2$ parallel to the $x$-axis, with the $y$-axis invariant, has matrix
\[
\begin{pmatrix}2&0\\0&1\end{pmatrix}
\]

Since the reflection is followed by the stretch,
\[
\mathbf Q=
\begin{pmatrix}2&0\\0&1\end{pmatrix}
\begin{pmatrix}0&1\\1&0\end{pmatrix}
=
\begin{pmatrix}0&2\\1&0\end{pmatrix}
\]

So
\[
\mathbf Q=\begin{pmatrix}0&2\\1&0\end{pmatrix}
\]

  \item From part (a), $p=2$, so
\[
\mathbf P=\begin{pmatrix}3&2\\2&3\end{pmatrix}
\]

Since $U$ followed by $V$ is represented by $\mathbf R$,
\[
\mathbf R=\mathbf Q\mathbf P
\]

Therefore
\begin{align*}
\mathbf R
&=
\begin{pmatrix}0&2\\1&0\end{pmatrix}
\begin{pmatrix}3&2\\2&3\end{pmatrix}\\
&=
\begin{pmatrix}
0\cdot 3+2\cdot 2 & 0\cdot 2+2\cdot 3\\
1\cdot 3+0\cdot 2 & 1\cdot 2+0\cdot 3
\end{pmatrix}\\
&=
\begin{pmatrix}4&6\\3&2\end{pmatrix}
\end{align*}

So
\[
\mathbf R=\begin{pmatrix}4&6\\3&2\end{pmatrix}
\]
\end{enumerate}
\end{tcolorbox}


\newpage

% ---- Question 3 ----
\questionitem

\begin{enumerate}[label=\textbf{(\alph*)}]
  \item Specify fully the transformation $\mathbf{T}$ of the plane associated with the matrix $\mathbf{M}$, where
  \[
  \mathbf{M} = \begin{pmatrix} \cos\theta & -\sin\theta \\ \sin\theta & \cos\theta \end{pmatrix}
  \]
  and $0<\theta<2\pi$. \marks{2} \\
  \item
  \begin{enumerate}[label=(\roman*)]
    \item Find $\det \mathbf{M}$. \marks{1}
    \item Deduce \textbf{two properties} of the transformation $\mathbf{T}$ from the value of $\det \mathbf{M}$. \marks{2} \\
  \end{enumerate}
  \item Prove that
  \[
  \mathbf{M}^n = \begin{pmatrix} \cos n\theta & -\sin n\theta \\ \sin n\theta & \cos n\theta \end{pmatrix}
  \]
  where $n$ is a positive integer. \marks{4} \\

  \item Hence specify fully a \textbf{single} transformation which is equivalent to $n$ applications of the transformation $\mathbf{T}$. \marks{1} \\
\end{enumerate}

\vspace*{30pt}

\begin{tcolorbox}[colback=white, colframe=scarlet, title={\textbf{Solution}}, breakable, grow to left by=6mm, grow to right by=10mm]
\begin{enumerate}[label=\textbf{(\alph*)}]
  \item For a point with position vector $\begin{pmatrix}x\\y\end{pmatrix}$, the transformation gives
  \[
  \begin{pmatrix}x'\\y'\end{pmatrix}
  =
  \begin{pmatrix}
  \cos\theta & -\sin\theta\\
  \sin\theta & \cos\theta
  \end{pmatrix}
  \begin{pmatrix}x\\y\end{pmatrix}
  \]

  This is the standard matrix for a rotation in the plane.

  So $\mathbf{T}$ is a rotation through angle $\theta$ anticlockwise about the origin.

  \item
  \begin{enumerate}[label=(\roman*)]
    \item
    \[
    \det \mathbf{M}
    =
    \begin{vmatrix}
    \cos\theta & -\sin\theta\\
    \sin\theta & \cos\theta
    \end{vmatrix}
    =
    \cos\theta\cos\theta-(-\sin\theta)(\sin\theta)
    \]
    \[
    =\cos^2\theta+\sin^2\theta=1
    \]

    Hence $\det \mathbf{M}=1$.

    \item Since $|\det \mathbf{M}|=1$, the area scale factor is $1$.

    So areas are preserved.

    Also, $\det \mathbf{M}>0$, so orientation is preserved.

    Therefore there is no reflection.

    The two properties are: area is preserved, and orientation is preserved.
  \end{enumerate}

  \item We prove the result by induction on $n$.

  First, when $n=1$,
  \[
  \mathbf{M}^1=\mathbf{M}=
  \begin{pmatrix}
  \cos\theta & -\sin\theta\\
  \sin\theta & \cos\theta
  \end{pmatrix}
  =
  \begin{pmatrix}
  \cos(1\theta) & -\sin(1\theta)\\
  \sin(1\theta) & \cos(1\theta)
  \end{pmatrix}
  \]

  So the statement is true for $n=1$.

  Now assume that for some positive integer $k$,
  \[
  \mathbf{M}^k=
  \begin{pmatrix}
  \cos(k\theta) & -\sin(k\theta)\\
  \sin(k\theta) & \cos(k\theta)
  \end{pmatrix}
  \]

  Then
  \[
  \mathbf{M}^{k+1}=\mathbf{M}^k\mathbf{M}
  \]

  Substituting the induction hypothesis,
  \[
  \mathbf{M}^{k+1}
  =
  \begin{pmatrix}
  \cos(k\theta) & -\sin(k\theta)\\
  \sin(k\theta) & \cos(k\theta)
  \end{pmatrix}
  \begin{pmatrix}
  \cos\theta & -\sin\theta\\
  \sin\theta & \cos\theta
  \end{pmatrix}
  \]

  Multiplying the matrices,
  \[
  \mathbf{M}^{k+1}
  =
  \begin{pmatrix}
  \cos(k\theta)\cos\theta-\sin(k\theta)\sin\theta
  &
  -\cos(k\theta)\sin\theta-\sin(k\theta)\cos\theta
  \\[4pt]
  \sin(k\theta)\cos\theta+\cos(k\theta)\sin\theta
  &
  -\sin(k\theta)\sin\theta+\cos(k\theta)\cos\theta
  \end{pmatrix}
  \]

  Using the angle addition formulae,
  \[
  \cos(A+B)=\cos A\cos B-\sin A\sin B
  \]
  \[
  \sin(A+B)=\sin A\cos B+\cos A\sin B
  \]

  we get
  \[
  \mathbf{M}^{k+1}
  =
  \begin{pmatrix}
  \cos((k+1)\theta) & -\sin((k+1)\theta)\\
  \sin((k+1)\theta) & \cos((k+1)\theta)
  \end{pmatrix}
  \]

  So the statement is true for $k+1$.

  Therefore, by induction,
  \[
  \mathbf{M}^n=
  \begin{pmatrix}
  \cos(n\theta) & -\sin(n\theta)\\
  \sin(n\theta) & \cos(n\theta)
  \end{pmatrix}
  \]
  for all positive integers $n$.

  \item From part (c), $n$ applications of $\mathbf{T}$ have matrix
  \[
  \mathbf{M}^n=
  \begin{pmatrix}
  \cos(n\theta) & -\sin(n\theta)\\
  \sin(n\theta) & \cos(n\theta)
  \end{pmatrix}
  \]

  This is the matrix of a rotation through angle $n\theta$ anticlockwise about the origin.

  So the single equivalent transformation is a rotation about the origin through angle $n\theta$ anticlockwise.
\end{enumerate}
\end{tcolorbox}


\newpage

% ---- Question 4 ----
\questionitem

Consider the matrix
\[
A=\begin{pmatrix}
-\frac{1}{2} & \frac{\sqrt{3}}{2}\\[6pt]
\frac{\sqrt{3}}{2} & \frac{1}{2}
\end{pmatrix}
\] \\

\begin{enumerate}[label=\textbf{(\roman*)}]
  \item Calculate $A^2$. \marks{1} \\
\end{enumerate}

You are told that the matrix $A$ represents a reflection in a line through the origin. \\

\begin{enumerate}[label=\textbf{(\roman*)}, resume]
  \item Explain why your answer to part (i) is consistent with this information. \marks{1} \\
  \item By finding the invariant points of the transformation, determine the equation of the mirror line. \marks{3} \\
  \item Describe fully the transformation represented by
  \[
  B=\begin{pmatrix}
  -\frac{1}{2} & -\frac{\sqrt{3}}{2}\\[6pt]
  \frac{\sqrt{3}}{2} & -\frac{1}{2}
  \end{pmatrix}
  \]
  \marks[-30pt]{2}
  \item A composite transformation is formed by the transformation represented by $A$ followed by the transformation represented by $B$. Find the single matrix that represents this composite transformation. \marks{2} \\
  \item The composite transformation described in part (v) is equivalent to a single reflection. Find the equation of the mirror line of this reflection. \marks{1} \\
\end{enumerate}

\vspace*{30pt}

\begin{tcolorbox}[colback=white, colframe=scarlet, title={\textbf{Solution}}, breakable, grow to left by=6mm, grow to right by=10mm]
\begin{enumerate}[label=\textbf{(\roman*)}]
\item
\[
A^2
=
\begin{pmatrix}
-\frac12 & \frac{\sqrt3}{2}\\[4pt]
\frac{\sqrt3}{2} & \frac12
\end{pmatrix}
\begin{pmatrix}
-\frac12 & \frac{\sqrt3}{2}\\[4pt]
\frac{\sqrt3}{2} & \frac12
\end{pmatrix}
\]

Calculating each entry,
\[
A^2
=
\begin{pmatrix}
\left(-\frac12\right)\left(-\frac12\right)+\left(\frac{\sqrt3}{2}\right)\left(\frac{\sqrt3}{2}\right)
&
\left(-\frac12\right)\left(\frac{\sqrt3}{2}\right)+\left(\frac{\sqrt3}{2}\right)\left(\frac12\right)
\\[8pt]
\left(\frac{\sqrt3}{2}\right)\left(-\frac12\right)+\left(\frac12\right)\left(\frac{\sqrt3}{2}\right)
&
\left(\frac{\sqrt3}{2}\right)\left(\frac{\sqrt3}{2}\right)+\left(\frac12\right)\left(\frac12\right)
\end{pmatrix}
\]

\[
=
\begin{pmatrix}
\frac14+\frac34 & -\frac{\sqrt3}{4}+\frac{\sqrt3}{4}\\[4pt]
-\frac{\sqrt3}{4}+\frac{\sqrt3}{4} & \frac34+\frac14
\end{pmatrix}
=
\begin{pmatrix}
1&0\\[2pt]
0&1
\end{pmatrix}
\]

So
\[
A^2=\begin{pmatrix}1&0\\0&1\end{pmatrix}
\]

\item
A reflection done twice takes every point back to its original position, so its matrix squared should be the identity matrix.

Since in part (i) we found \(A^2=I\), this is consistent with \(A\) representing a reflection.

\item
Points on the mirror line are unchanged by the reflection, so invariant points satisfy
\[
A\begin{pmatrix}x\\y\end{pmatrix}
=
\begin{pmatrix}x\\y\end{pmatrix}
\]

Thus
\[
\begin{pmatrix}
-\frac12 & \frac{\sqrt3}{2}\\[4pt]
\frac{\sqrt3}{2} & \frac12
\end{pmatrix}
\begin{pmatrix}x\\y\end{pmatrix}
=
\begin{pmatrix}x\\y\end{pmatrix}
\]

This gives the simultaneous equations
\begin{align*}
-\frac12x+\frac{\sqrt3}{2}y&=x\\
\frac{\sqrt3}{2}x+\frac12y&=y
\end{align*}

Multiply both equations by \(2\):
\begin{align*}
-x+\sqrt3\,y&=2x\\
\sqrt3\,x+y&=2y
\end{align*}

So
\begin{align*}
-3x+\sqrt3\,y&=0\\
\sqrt3\,x-y&=0
\end{align*}

From either equation,
\[
y=\sqrt3\,x
\]

So the mirror line is
\[
y=\sqrt3 x
\]

\item
The standard matrix for a rotation through angle \(\theta\) about the origin is
\[
\begin{pmatrix}
\cos\theta & -\sin\theta\\
\sin\theta & \cos\theta
\end{pmatrix}
\]

Now
\[
B=
\begin{pmatrix}
-\frac12 & -\frac{\sqrt3}{2}\\[4pt]
\frac{\sqrt3}{2} & -\frac12
\end{pmatrix}
\]

So
\[
\cos\theta=-\frac12,\qquad \sin\theta=\frac{\sqrt3}{2}
\]

This gives
\[
\theta=120^\circ
\]

Therefore \(B\) represents a rotation of \(120^\circ\) anticlockwise about the origin.

\item
If \(A\) is followed by \(B\), the composite matrix is \(BA\).

\[
BA=
\begin{pmatrix}
-\frac12 & -\frac{\sqrt3}{2}\\[4pt]
\frac{\sqrt3}{2} & -\frac12
\end{pmatrix}
\begin{pmatrix}
-\frac12 & \frac{\sqrt3}{2}\\[4pt]
\frac{\sqrt3}{2} & \frac12
\end{pmatrix}
\]

\[
=
\begin{pmatrix}
\left(-\frac12\right)\left(-\frac12\right)+\left(-\frac{\sqrt3}{2}\right)\left(\frac{\sqrt3}{2}\right)
&
\left(-\frac12\right)\left(\frac{\sqrt3}{2}\right)+\left(-\frac{\sqrt3}{2}\right)\left(\frac12\right)
\\[8pt]
\left(\frac{\sqrt3}{2}\right)\left(-\frac12\right)+\left(-\frac12\right)\left(\frac{\sqrt3}{2}\right)
&
\left(\frac{\sqrt3}{2}\right)\left(\frac{\sqrt3}{2}\right)+\left(-\frac12\right)\left(\frac12\right)
\end{pmatrix}
\]

\[
=
\begin{pmatrix}
\frac14-\frac34 & -\frac{\sqrt3}{4}-\frac{\sqrt3}{4}\\[4pt]
-\frac{\sqrt3}{4}-\frac{\sqrt3}{4} & \frac34-\frac14
\end{pmatrix}
=
\begin{pmatrix}
-\frac12 & -\frac{\sqrt3}{2}\\[4pt]
-\frac{\sqrt3}{2} & \frac12
\end{pmatrix}
\]

So the single matrix is
\[
\begin{pmatrix}
-\frac12 & -\frac{\sqrt3}{2}\\[4pt]
-\frac{\sqrt3}{2} & \frac12
\end{pmatrix}
\]

\item
Let
\[
M=
\begin{pmatrix}
-\frac12 & -\frac{\sqrt3}{2}\\[4pt]
-\frac{\sqrt3}{2} & \frac12
\end{pmatrix}
\]

The mirror line consists of invariant points, so solve
\[
M\begin{pmatrix}x\\y\end{pmatrix}
=
\begin{pmatrix}x\\y\end{pmatrix}
\]

This gives
\begin{align*}
-\frac12x-\frac{\sqrt3}{2}y&=x\\
-\frac{\sqrt3}{2}x+\frac12y&=y
\end{align*}

Multiplying by \(2\),
\begin{align*}
-x-\sqrt3\,y&=2x\\
-\sqrt3\,x+y&=2y
\end{align*}

So
\begin{align*}
-3x-\sqrt3\,y&=0\\
-\sqrt3\,x-y&=0
\end{align*}

Hence
\[
y=-\sqrt3 x
\]

So the mirror line is
\[
y=-\sqrt3 x
\]
\end{enumerate}
\end{tcolorbox}


\newpage

% ---- Question 5 ----
\questionitem

The transformation $U$, represented by the $2\times 2$ matrix $P$, is a reflection in the $x$-axis. \\

\begin{enumerate}[label=\textbf{(\alph*)}]
  \item Write down the matrix $P$. \marks{1} \\
\end{enumerate}

The transformation $V$, represented by the $2\times 2$ matrix $Q$, is a reflection in the line $y=x$. \\

\begin{enumerate}[label=\textbf{(\alph*)}, resume]
  \item Write down the matrix $Q$. \marks{1} \\
\end{enumerate}

Given that $U$ followed by $V$ is transformation $T$, which is represented by the matrix $R$ \\

\begin{enumerate}[label=\textbf{(\alph*)}, resume]
  \item express $R$ in terms of $P$ and $Q$, \marks{1}
  \item find the matrix $R$, \marks{2}
  \item give a full geometrical description of $T$ as a single transformation. \marks{2} \\
\end{enumerate}

\vspace*{30pt}

\begin{tcolorbox}[colback=white, colframe=scarlet, title={\textbf{Solution}}, breakable, grow to left by=6mm, grow to right by=10mm]
\begin{enumerate}[label=\textbf{(\alph*)}]
  \item A reflection in the $x$-axis maps $(x,y)$ to $(x,-y)$, so
  \[
  P=\begin{pmatrix}1&0\\0&-1\end{pmatrix}
  \]

  \item A reflection in the line $y=x$ swaps the coordinates, so
  \[
  Q=\begin{pmatrix}0&1\\1&0\end{pmatrix}
  \]

  \item If a point is first transformed by $U$ and then by $V$, then with column vectors
  \[
  \mathbf{x}\mapsto P\mathbf{x}\mapsto Q(P\mathbf{x})=(QP)\mathbf{x}
  \]
  Hence
  \[
  R=QP
  \]

  \item Using part (c),
  \[
  R=QP=\begin{pmatrix}0&1\\1&0\end{pmatrix}\begin{pmatrix}1&0\\0&-1\end{pmatrix}
  \]
  Multiplying,
  \[
  R=\begin{pmatrix}
  0\cdot 1+1\cdot 0 & 0\cdot 0+1\cdot (-1)\\
  1\cdot 1+0\cdot 0 & 1\cdot 0+0\cdot (-1)
  \end{pmatrix}
  =\begin{pmatrix}0&-1\\1&0\end{pmatrix}
  \]
  So the matrix $R$ is
  \[
  \begin{pmatrix}0&-1\\1&0\end{pmatrix}
  \]

  \item Apply $R$ to a general point:
  \[
  \begin{pmatrix}x'\\y'\end{pmatrix}
  =
  \begin{pmatrix}0&-1\\1&0\end{pmatrix}
  \begin{pmatrix}x\\y\end{pmatrix}
  =
  \begin{pmatrix}-y\\x\end{pmatrix}
  \]
  So $(x,y)\mapsto (-y,x)$.

  This is a rotation through $90^\circ$ anticlockwise about the origin.

  Therefore $T$ is a rotation of $90^\circ$ anticlockwise about the origin.
\end{enumerate}
\end{tcolorbox}


\newpage

% ---- Question 6 ----
\questionitem

The matrix $\mathbf{M}$ is given by
\[
\mathbf{M} = \begin{pmatrix} 0 & 1 \\ 1 & 0 \end{pmatrix}\begin{pmatrix} 1 & 0 \\ 0 & k \end{pmatrix}
\]
where $k>0$ and $k \ne 1$. \\

\begin{enumerate}[label=\textbf{(\alph*)}]
  \item The matrix $\mathbf{M}$ represents a sequence of two geometrical transformations. State the type of each transformation, and make clear the order in which they are applied. \marks{2} \\
\end{enumerate}

The unit square in the $x$-$y$ plane is transformed by $\mathbf{M}$ onto rectangle $OABC$. \\

\begin{enumerate}[label=\textbf{(\alph*)}, resume]
  \item Find, in terms of $k$, the area of rectangle $OABC$ and the matrix which transforms $OABC$ onto the unit square. \marks{3} \\
  \item Show that the line through the origin with gradient $\frac{1}{\sqrt{k}}$ is invariant under the transformation in the $x$-$y$ plane represented by $\mathbf{M}$. \marks{3} \\
\end{enumerate}

\vspace*{30pt}

\begin{tcolorbox}[colback=white, colframe=scarlet, title={\textbf{Solution}}, breakable, grow to left by=6mm, grow to right by=10mm]
\begin{enumerate}[label=\textbf{(\alph*)}]
  \item Write
  \[
  \mathbf{M}
  =
  \begin{pmatrix}
  0 & 1\\
  1 & 0
  \end{pmatrix}
  \begin{pmatrix}
  1 & 0\\
  0 & k
  \end{pmatrix}
  \]
  
  The right-hand matrix is applied first. Since
  \[
  \begin{pmatrix}
  1 & 0\\
  0 & k
  \end{pmatrix}
  \begin{pmatrix}
  x\\y
  \end{pmatrix}
  =
  \begin{pmatrix}
  x\\ky
  \end{pmatrix},
  \]
  this is a stretch parallel to the $y$-axis, with scale factor $k$.

  Then
  \[
  \begin{pmatrix}
  0 & 1\\
  1 & 0
  \end{pmatrix}
  \begin{pmatrix}
  x\\y
  \end{pmatrix}
  =
  \begin{pmatrix}
  y\\x
  \end{pmatrix},
  \]
  which swaps the coordinates, so this is a reflection in the line $y=x$.

  Therefore the transformations are: first a stretch in the $y$-direction by scale factor $k$, then a reflection in the line $y=x$.

  \item First find $\mathbf{M}$:
  \[
  \mathbf{M}
  =
  \begin{pmatrix}
  0 & 1\\
  1 & 0
  \end{pmatrix}
  \begin{pmatrix}
  1 & 0\\
  0 & k
  \end{pmatrix}
  =
  \begin{pmatrix}
  0 & k\\
  1 & 0
  \end{pmatrix}
  \]

  The area scale factor of a transformation is $|\det \mathbf{M}|$, so
  \[
  \det \mathbf{M}
  =
  \begin{vmatrix}
  0 & k\\
  1 & 0
  \end{vmatrix}
  =0-k=-k
  \]
  Hence the area scale factor is
  \[
  |\det \mathbf{M}|=k
  \]
  Since the unit square has area $1$, the area of rectangle $OABC$ is
  \[
  1\times k=k
  \]

  To transform $OABC$ back to the unit square, we need $\mathbf{M}^{-1}$.
  For
  \[
  \mathbf{M}=
  \begin{pmatrix}
  0 & k\\
  1 & 0
  \end{pmatrix},
  \]
  \[
  \mathbf{M}^{-1}
  =
  \frac{1}{-k}
  \begin{pmatrix}
  0 & -k\\
  -1 & 0
  \end{pmatrix}
  =
  \begin{pmatrix}
  0 & 1\\
  \frac{1}{k} & 0
  \end{pmatrix}
  \]

  Therefore the area of $OABC$ is $k$, and the required matrix is
  \[
  \begin{pmatrix}
  0 & 1\\
  \frac{1}{k} & 0
  \end{pmatrix}
  \]

  \item The line through the origin with gradient $\dfrac{1}{\sqrt{k}}$ has equation
  \[
  y=\frac{x}{\sqrt{k}}
  \]

  Take any point on this line:
  \[
  \begin{pmatrix}
  x\\[2pt]
  \dfrac{x}{\sqrt{k}}
  \end{pmatrix}
  \]

  Apply $\mathbf{M}=\begin{pmatrix}0&k\\1&0\end{pmatrix}$:
  \[
  \mathbf{M}
  \begin{pmatrix}
  x\\[2pt]
  \dfrac{x}{\sqrt{k}}
  \end{pmatrix}
  =
  \begin{pmatrix}
  0 & k\\
  1 & 0
  \end{pmatrix}
  \begin{pmatrix}
  x\\[2pt]
  \dfrac{x}{\sqrt{k}}
  \end{pmatrix}
  =
  \begin{pmatrix}
  \dfrac{kx}{\sqrt{k}}\\[4pt]
  x
  \end{pmatrix}
  =
  \begin{pmatrix}
  \sqrt{k}\,x\\
  x
  \end{pmatrix}
  \]

  Let the image point be $(X,Y)=(\sqrt{k}\,x,\ x)$. Then
  \[
  \frac{X}{\sqrt{k}}=\frac{\sqrt{k}\,x}{\sqrt{k}}=x=Y
  \]
  so
  \[
  Y=\frac{X}{\sqrt{k}}
  \]

  Therefore the image point also lies on the line $y=\dfrac{x}{\sqrt{k}}$.

  Hence the line through the origin with gradient $\dfrac{1}{\sqrt{k}}$ is invariant under the transformation represented by $\mathbf{M}$.
\end{enumerate}
\end{tcolorbox}


\newpage

% ---- Question 7 ----
\questionitem

You are given the matrix $\mathbf{A} = \begin{pmatrix}0 & 1 & 0\\0 & 0 & 1\\1 & 0 & 0\end{pmatrix}$. \\

\begin{enumerate}[label=\textbf{(\alph*)}]
  \item Find $\mathbf{A}^5$. \marks{1} \\
  \item Describe the transformation that $\mathbf{A}$ represents. \marks{2} \\
\end{enumerate}

The matrix $\mathbf{B}$ represents a reflection in the plane $x = y$. \\

\begin{enumerate}[label=\textbf{(\alph*)}, resume]
  \item Write down the matrix $\mathbf{B}$. \marks{1} \\
\end{enumerate}

The point $Q$ has coordinates $(4, -1, 2)$. The point $Q'$ is the image of $Q$ under the transformation represented by $\mathbf{B}$. \\

\begin{enumerate}[label=\textbf{(\alph*)}, resume]
  \item Find the coordinates of $Q'$. \marks{1} \\
\end{enumerate}

\vspace*{30pt}

\begin{tcolorbox}[colback=white, colframe=scarlet, title={\textbf{Solution}}, breakable, grow to left by=6mm, grow to right by=10mm]
\begin{enumerate}[label=\textbf{(\alph*)}]
  \item First find \(\mathbf{A}^2\):
\[
\mathbf{A}^2
=
\begin{pmatrix}
0 & 1 & 0\\
0 & 0 & 1\\
1 & 0 & 0
\end{pmatrix}
\begin{pmatrix}
0 & 1 & 0\\
0 & 0 & 1\\
1 & 0 & 0
\end{pmatrix}
=
\begin{pmatrix}
0 & 0 & 1\\
1 & 0 & 0\\
0 & 1 & 0
\end{pmatrix}
\]

Now multiply once more:
\[
\mathbf{A}^3
=
\mathbf{A}^2\mathbf{A}
=
\begin{pmatrix}
0 & 0 & 1\\
1 & 0 & 0\\
0 & 1 & 0
\end{pmatrix}
\begin{pmatrix}
0 & 1 & 0\\
0 & 0 & 1\\
1 & 0 & 0
\end{pmatrix}
=
\begin{pmatrix}
1 & 0 & 0\\
0 & 1 & 0\\
0 & 0 & 1
\end{pmatrix}
=
\mathbf{I}
\]

So
\[
\mathbf{A}^5=\mathbf{A}^3\mathbf{A}^2=\mathbf{I}\mathbf{A}^2=\mathbf{A}^2
\]

Hence
\[
\mathbf{A}^5=
\begin{pmatrix}
0 & 0 & 1\\
1 & 0 & 0\\
0 & 1 & 0
\end{pmatrix}
\]

  \item Let a general point be \(\begin{pmatrix}x\\y\\z\end{pmatrix}\). Then
\[
\mathbf{A}\begin{pmatrix}x\\y\\z\end{pmatrix}
=
\begin{pmatrix}
0 & 1 & 0\\
0 & 0 & 1\\
1 & 0 & 0
\end{pmatrix}
\begin{pmatrix}x\\y\\z\end{pmatrix}
=
\begin{pmatrix}y\\z\\x\end{pmatrix}
\]

So \(\mathbf{A}\) sends \((x,y,z)\) to \((y,z,x)\), which is a cyclic permutation of the coordinates.

Also, any point on the line \(x=y=z\) is unchanged, so this is a rotation about the line \(x=y=z\).

Since \(\mathbf{A}^3=\mathbf{I}\), three applications give a full turn, so the angle of rotation is
\[
\frac{360^\circ}{3}=120^\circ
\]

Therefore \(\mathbf{A}\) represents a rotation of \(120^\circ\) about the line \(x=y=z\).

  \item A reflection in the plane \(x=y\) swaps the \(x\)- and \(y\)-coordinates and leaves \(z\) unchanged.

Therefore
\[
\mathbf{B}=
\begin{pmatrix}
0 & 1 & 0\\
1 & 0 & 0\\
0 & 0 & 1
\end{pmatrix}
\]

  \item Apply \(\mathbf{B}\) to \(Q=(4,-1,2)\):
\[
\mathbf{Q}'
=
\mathbf{B}
\begin{pmatrix}
4\\
-1\\
2
\end{pmatrix}
=
\begin{pmatrix}
0 & 1 & 0\\
1 & 0 & 0\\
0 & 0 & 1
\end{pmatrix}
\begin{pmatrix}
4\\
-1\\
2
\end{pmatrix}
=
\begin{pmatrix}
-1\\
4\\
2
\end{pmatrix}
\]

So the coordinates of \(Q'\) are \((-1,4,2)\).
\end{enumerate}
\end{tcolorbox}


\newpage

% ---- Question 8 ----
\questionitem

A right-angled triangle $T$ has vertices $A(-2,1)$, $B(1,1)$ and $C(1,5)$. When $T$ is transformed by the matrix
\[
P = \begin{pmatrix} 0 & -1 \\ 1 & 0 \end{pmatrix}
\]
the image is $T'$. \\

\begin{enumerate}[label=\textbf{(\alph*)}]
  \item Find the coordinates of the vertices of $T'$. \marks{2} \\
  \item Describe fully the transformation represented by $P$. \marks{2} \\
\end{enumerate}

The matrices
\[
Q = \begin{pmatrix} 2 & -1 \\ 0 & 3 \end{pmatrix}
\]
and
\[
R = \begin{pmatrix} 1 & 2 \\ 1 & 0 \end{pmatrix}
\]
represent two transformations. When $T$ is transformed by the matrix $QR$, the image is $T''$. \\

\begin{enumerate}[label=\textbf{(\alph*)}, resume]
  \item Find $QR$. \marks{2} \\
  \item Find the determinant of $QR$. \marks{2} \\
  \item Using your answer to part (d), find the area of $T''$. \marks{3} \\
\end{enumerate}

\vspace*{30pt}

\begin{tcolorbox}[colback=white, colframe=scarlet, title={\textbf{Solution}}, breakable, grow to left by=6mm, grow to right by=10mm]
\begin{enumerate}[label=\textbf{(\alph*)}]
  \item A point \((x,y)\) is represented by the column vector \(\begin{pmatrix}x\\y\end{pmatrix}\).

  Applying
  \[
  P=\begin{pmatrix}0&-1\\1&0\end{pmatrix}
  \]
  gives
  \[
  P\begin{pmatrix}x\\y\end{pmatrix}
  =
  \begin{pmatrix}0&-1\\1&0\end{pmatrix}\begin{pmatrix}x\\y\end{pmatrix}
  =
  \begin{pmatrix}-y\\x\end{pmatrix}
  \]

  So each point \((x,y)\) maps to \((-y,x)\).

  For \(A(-2,1)\),
  \[
  A' = (-1,-2)
  \]

  For \(B(1,1)\),
  \[
  B' = (-1,1)
  \]

  For \(C(1,5)\),
  \[
  C' = (-5,1)
  \]

  Hence the vertices of \(T'\) are \(A'(-1,-2)\), \(B'(-1,1)\), \(C'(-5,1)\).

  \item To identify the transformation, look at what \(P\) does to the unit basis vectors.

  \[
  P\begin{pmatrix}1\\0\end{pmatrix}
  =
  \begin{pmatrix}0\\1\end{pmatrix}
  \]
  so the positive \(x\)-axis maps to the positive \(y\)-axis.

  Also,
  \[
  P\begin{pmatrix}0\\1\end{pmatrix}
  =
  \begin{pmatrix}-1\\0\end{pmatrix}
  \]
  so the positive \(y\)-axis maps to the negative \(x\)-axis.

  This is a rotation of \(90^\circ\) anticlockwise about the origin.

  \item
  \[
  QR=
  \begin{pmatrix}2&-1\\0&3\end{pmatrix}
  \begin{pmatrix}1&2\\1&0\end{pmatrix}
  \]

  Multiplying row by column:
  \[
  QR=
  \begin{pmatrix}
  2(1)+(-1)(1) & 2(2)+(-1)(0)\\
  0(1)+3(1) & 0(2)+3(0)
  \end{pmatrix}
  =
  \begin{pmatrix}
  1&4\\
  3&0
  \end{pmatrix}
  \]

  So
  \[
  QR=\begin{pmatrix}1&4\\3&0\end{pmatrix}
  \]

  \item Using
  \[
  QR=\begin{pmatrix}1&4\\3&0\end{pmatrix}
  \]
  its determinant is
  \[
  \det(QR)=1\cdot 0-4\cdot 3=-12
  \]

  Therefore,
  \[
  \det(QR)=-12
  \]

  \item The original triangle \(T\) is right-angled with side lengths
  \[
  AB=3,\qquad BC=4
  \]
  so its area is
  \[
  \text{Area of }T=\frac12\times 3\times 4=6
  \]

  The magnitude of the determinant gives the area scale factor, so
  \[
  \text{area scale factor}=|\det(QR)|=|-12|=12
  \]

  Therefore,
  \[
  \text{Area of }T''=12\times 6=72
  \]

  So the area of \(T''\) is \(72\).
\end{enumerate}
\end{tcolorbox}


\newpage

% ---- Question 9 ----
\questionitem

Given that the matrix $\mathbf{M}$, where
\[
\mathbf{M} = \begin{pmatrix}
-\frac{\sqrt{2}}{2} & 0 & -\frac{\sqrt{2}}{2} \\
0 & 1 & 0 \\
\frac{\sqrt{2}}{2} & 0 & -\frac{\sqrt{2}}{2}
\end{pmatrix}
\]
represents a rotation through $\alpha^\circ$ about the $y$-axis in the positive sense given by the right-hand rule. \\

\begin{enumerate}[label=\textbf{(\alph*)}]
  \item determine the value of $\alpha$. \marks{1} \\
  \item Hence determine the smallest possible positive integer value of $k$ for which $\mathbf{M}^k = \mathbf{I}$. \marks{2} \\
\end{enumerate}

The $3\times 3$ matrix $\mathbf{N}$ represents a reflection in the plane with equation $x = 0$. \\

\begin{enumerate}[label=\textbf{(\alph*)}, resume]
  \item Write down the matrix $\mathbf{N}$. \marks{1} \\
\end{enumerate}

The point $A$ has coordinates $(1, 2, 3)$. \\

The point $B$ is the image of the point $A$ under the transformation represented by matrix $\mathbf{M}$ followed by the transformation represented by matrix $\mathbf{N}$. \\

\begin{enumerate}[label=\textbf{(\alph*)}, resume]
  \item Show that the coordinates of $B$ are $(2\sqrt{2}, 2, -\sqrt{2})$. \marks{2} \\
\end{enumerate}

Given that $O$ is the origin \\

\begin{enumerate}[label=\textbf{(\alph*)}, resume]
  \item show that, to $3$ significant figures, the size of angle $AOB$ is $79.4^\circ$. \marks{2} \\
  \item Hence determine the area of triangle $AOB$, giving your answer to $3$ significant figures. \marks{2} \\
\end{enumerate}

\vspace*{30pt}

\begin{tcolorbox}[colback=white, colframe=scarlet, title={\textbf{Solution}}, breakable, grow to left by=6mm, grow to right by=10mm]
\begin{enumerate}[label=\textbf{(\alph*)}]
\item Comparing \(\mathbf M\) with the standard rotation matrix about the \(y\)-axis,
\[
\begin{pmatrix}
\cos\alpha & 0 & \sin\alpha\\
0&1&0\\
-\sin\alpha&0&\cos\alpha
\end{pmatrix}
\]
gives
\[
\cos\alpha=-\frac{\sqrt2}{2},
\qquad
\sin\alpha=-\frac{\sqrt2}{2}
\]
So \(\alpha\) is in the third quadrant, hence
\[
\alpha=225^\circ
\]

\item \(\mathbf M\) represents a rotation of \(225^\circ\), so \(\mathbf M^k\) represents a rotation of \(225k^\circ\).

For \(\mathbf M^k=\mathbf I\), the total rotation must be a multiple of \(360^\circ\), so
\[
225k=360n
\]
for some positive integer \(n\).

Dividing by \(45\),
\[
5k=8n
\]
The smallest positive integer \(k\) satisfying this is \(k=8\) (with \(n=5\)).

So the smallest possible value of \(k\) is \(8\).

\item Reflection in the plane \(x=0\) changes \((x,y,z)\) to \((-x,y,z)\), so
\[
\mathbf N=
\begin{pmatrix}
-1&0&0\\
0&1&0\\
0&0&1
\end{pmatrix}
\]

\item The point \(A\) has position vector
\[
\mathbf a=
\begin{pmatrix}
1\\
2\\
3
\end{pmatrix}
\]
First apply \(\mathbf M\):
\begin{align*}
\mathbf M\mathbf a
&=
\begin{pmatrix}
-\frac{\sqrt2}{2} & 0 & -\frac{\sqrt2}{2} \\
0 & 1 & 0 \\
\frac{\sqrt2}{2} & 0 & -\frac{\sqrt2}{2}
\end{pmatrix}
\begin{pmatrix}
1\\
2\\
3
\end{pmatrix} \\
&=
\begin{pmatrix}
-\frac{\sqrt2}{2}(1)-\frac{\sqrt2}{2}(3)\\
2\\
\frac{\sqrt2}{2}(1)-\frac{\sqrt2}{2}(3)
\end{pmatrix} \\
&=
\begin{pmatrix}
-\frac{4\sqrt2}{2}\\
2\\
-\frac{2\sqrt2}{2}
\end{pmatrix}
=
\begin{pmatrix}
-2\sqrt2\\
2\\
-\sqrt2
\end{pmatrix}
\end{align*}
Then apply \(\mathbf N\):
\begin{align*}
\mathbf N(\mathbf M\mathbf a)
&=
\begin{pmatrix}
-1&0&0\\
0&1&0\\
0&0&1
\end{pmatrix}
\begin{pmatrix}
-2\sqrt2\\
2\\
-\sqrt2
\end{pmatrix} \\
&=
\begin{pmatrix}
2\sqrt2\\
2\\
-\sqrt2
\end{pmatrix}
\end{align*}
Hence the coordinates of \(B\) are
\[
(2\sqrt2,\,2,\,-\sqrt2)
\]

\item \(\overrightarrow{OA}=(1,2,3)\) and \(\overrightarrow{OB}=(2\sqrt2,2,-\sqrt2)\).

Using the scalar product,
\[
\overrightarrow{OA}\cdot\overrightarrow{OB}
=1(2\sqrt2)+2(2)+3(-\sqrt2)
=4-\sqrt2
\]
Also,
\[
|\overrightarrow{OA}|=\sqrt{1^2+2^2+3^2}=\sqrt{14}
\]
and
\[
|\overrightarrow{OB}|=\sqrt{(2\sqrt2)^2+2^2+(-\sqrt2)^2}
=\sqrt{8+4+2}
=\sqrt{14}
\]
So if \(\theta=\angle AOB\), then
\[
\cos\theta
=\frac{\overrightarrow{OA}\cdot\overrightarrow{OB}}
{|\overrightarrow{OA}|\,|\overrightarrow{OB}|}
=\frac{4-\sqrt2}{14}
\]
Therefore
\[
\theta=\cos^{-1}\!\left(\frac{4-\sqrt2}{14}\right)\approx 79.4^\circ
\]
So the size of angle \(AOB\) is \(79.4^\circ\) to \(3\) significant figures.

\item Using the area formula for a triangle,
\[
\text{Area of } \triangle AOB
=\frac12 |\overrightarrow{OA}|\,|\overrightarrow{OB}|\,\sin\angle AOB
\]
From part (e), \(|\overrightarrow{OA}|=|\overrightarrow{OB}|=\sqrt{14}\), so
\[
\text{Area}
=\frac12(\sqrt{14})(\sqrt{14})\sin\theta
=7\sin\theta
\]
Using
\[
\cos\theta=\frac{4-\sqrt2}{14}
\]
we get
\[
\sin\theta=\sqrt{1-\left(\frac{4-\sqrt2}{14}\right)^2}
\]
Hence
\begin{align*}
\text{Area}
&=7\sqrt{1-\left(\frac{4-\sqrt2}{14}\right)^2} \\
&\approx 6.88
\end{align*}
So the area of triangle \(AOB\) is \(6.88\) square units to \(3\) significant figures.
\end{enumerate}
\end{tcolorbox}


\newpage

% ---- Question 10 ----
\questionitem

Bilal has been asked to describe the transformation given by the matrix
\[
\begin{pmatrix}
\frac{\sqrt{3}}{2} & 0 & -\frac{1}{2} \\
0 & 1 & 0 \\
\frac{1}{2} & 0 & \frac{\sqrt{3}}{2}
\end{pmatrix}
\]

He writes his answer as follows: \\


\fbox{%
\parbox{\dimexpr\linewidth-2\fboxsep-2\fboxrule\relax}{%
The transformation is a rotation about the $y$-axis through an angle of $\theta$, where
\[
\cos \theta = \frac{\sqrt{3}}{2} \quad \text{and} \quad \sin \theta = \frac{1}{2}
\]
\[
\theta = 30^\circ
\]
}%
} \\[10pt]

Identify and correct the error in Bilal's work. \marks{2}

\vspace*{30pt}

\begin{tcolorbox}[colback=white, colframe=scarlet, title={\textbf{Solution}}, breakable, grow to left by=6mm, grow to right by=10mm]
The standard matrix for a rotation about the $y$-axis is
\[
R_y(\theta)=
\begin{bmatrix}
\cos\theta & 0 & \sin\theta\\
0 & 1 & 0\\
-\sin\theta & 0 & \cos\theta
\end{bmatrix}
\]

Comparing this with
\[
\begin{bmatrix}
\frac{\sqrt3}{2} & 0 & -\frac12\\
0 & 1 & 0\\
\frac12 & 0 & \frac{\sqrt3}{2}
\end{bmatrix}
\]
gives
\[
\cos\theta=\frac{\sqrt3}{2},\qquad \sin\theta=-\frac12
\]

So the angle is not $30^\circ$, but
\[
\theta=-30^\circ
\]
equivalently $330^\circ$.

Bilal's error is that he ignored the negative sign on the sine term. The transformation is a rotation about the $y$-axis through $30^\circ$ clockwise when viewed from the positive $y$-axis.
\end{tcolorbox}


\newpage

% ---- Question 11 ----
\questionitem

Three non-singular $3 \times 3$ matrices, $\mathbf{A}$, $\mathbf{B}$ and $\mathbf{R}$, are such that
\[
\mathbf{R}\mathbf{A}=\mathbf{B}
\] \\
The matrix $\mathbf{R}$ represents a rotation about the $x$-axis through an angle $\theta$ and
\[
\mathbf{B}=\begin{bmatrix}
1 & 0 & 0 \\
0 & -\sin\theta & \cos\theta \\
0 & \cos\theta & \sin\theta
\end{bmatrix}
\] \\
\begin{enumerate}[label=\textbf{(\alph*)}]
  \item Show that $\mathbf{A}$ is independent of the value of $\theta$. \marks{3} \\
  \item Give a full description of the single transformation represented by the matrix $\mathbf{A}$. \marks{1} \\
\end{enumerate}

\vspace*{30pt}

\begin{tcolorbox}[colback=white, colframe=scarlet, title={\textbf{Solution}}, breakable, grow to left by=6mm, grow to right by=10mm]
\begin{enumerate}[label=\textbf{(\alph*)}]
  \item Since
\[
\mathbf{R}\mathbf{A}=\mathbf{B}
\]
and $\mathbf{R}$ is non-singular, we can multiply on the left by $\mathbf{R}^{-1}$:
\[
\mathbf{A}=\mathbf{R}^{-1}\mathbf{B}
\]

For a rotation about the $x$-axis through angle $\theta$,
\[
\mathbf{R}=
\begin{bmatrix}
1&0&0\\
0&\cos\theta&-\sin\theta\\
0&\sin\theta&\cos\theta
\end{bmatrix}
\]
so
\[
\mathbf{R}^{-1}=
\begin{bmatrix}
1&0&0\\
0&\cos\theta&\sin\theta\\
0&-\sin\theta&\cos\theta
\end{bmatrix}
\]

Now calculate $\mathbf{A}$:
\begin{align*}
\mathbf{A}
&=
\begin{bmatrix}
1&0&0\\
0&\cos\theta&\sin\theta\\
0&-\sin\theta&\cos\theta
\end{bmatrix}
\begin{bmatrix}
1&0&0\\
0&-\sin\theta&\cos\theta\\
0&\cos\theta&\sin\theta
\end{bmatrix} \\[6pt]
&=
\begin{bmatrix}
1&0&0\\
0&-\cos\theta\sin\theta+\sin\theta\cos\theta&\cos^2\theta+\sin^2\theta\\
0&\sin^2\theta+\cos^2\theta&-\sin\theta\cos\theta+\cos\theta\sin\theta
\end{bmatrix} \\[6pt]
&=
\begin{bmatrix}
1&0&0\\
0&0&1\\
0&1&0
\end{bmatrix}
\end{align*}

This matrix contains no $\theta$, so $\mathbf{A}$ is independent of the value of $\theta$.

Hence
\[
\mathbf{A}=
\begin{bmatrix}
1&0&0\\
0&0&1\\
0&1&0
\end{bmatrix}
\]

  \item The matrix $\mathbf{A}$ maps
\[
\begin{bmatrix}x\\y\\z\end{bmatrix}
\mapsto
\begin{bmatrix}x\\z\\y\end{bmatrix}
\]
so it leaves $x$ unchanged and swaps $y$ and $z$.

Therefore, $\mathbf{A}$ represents a reflection in the plane $y=z$.
\end{enumerate}
\end{tcolorbox}


\end{enumerate}
\end{document}
